\chapter{CONSIDERAÇÕES FINAIS}

Este trabalho buscou investigar a viabilidade, os desafios, o desempenho comparativo e as estratégias de integração do protocolo HTTP/3 em APIs REST desenvolvidas com Node.js, conforme o objetivo estabelecido na Introdução. A fundamentação teórica (Capítulo 2) e a revisão rápida da literatura (Capítulo 3) forneceram o embasamento para o estudo de caso experimental descrito no Capítulo 4 e para a análise dos resultados apresentada no Capítulo 5. As seções a seguir sintetizam os achados em relação às questões de pesquisa, explicitam as contribuições do estudo, remetem às limitações já discutidas quantitativamente e indicam direções para pesquisas posteriores.

\section{Síntese dos resultados e respostas às questões de pesquisa}

Acerca da viabilidade do HTTP/3 em Node.js sem suporte nativo estável (QP1), os experimentos confirmam que é viável expor HTTP/3 a uma API Node.js que permanece servida por HTTP/1.1 no upstream, mediante terminação QUIC/TLS em um proxy reverso NGINX. Nas duas campanhas reportadas no Capítulo 5, o caminho HTTP/3 respondeu com taxa de falha nula em todas as células agregadas - 12 células (120 execuções) na submatriz degradada e 21 células (210 execuções) na matriz ideal ---, o que indica operação estável do fluxo cliente→NGINX→API nos ambientes construídos. Esse resultado está alinhado à literatura que discute a adoção de HTTP/3 em stacks ainda em consolidação e à especificação normativa de QUIC e HTTP/3 (KOCH; FALOWO; ELROD, 2024; QUIC Working Group - IETF, 2021; IETF, 2022).

Entre os principais desafios técnicos (QP2), a implementação e a medição revelaramse operacionalmente complexas: terminação TLS/QUIC na borda, salto adicional entre NGINX e a API, observabilidade fragmentada em camadas distintas, e assimetria entre ferramentas k6 para HTTP/1.1 e HTTP/2, cliente QUIC em contêiner para HTTP/3. Essa assimetria aumentou a variabilidade entre réplicas no caminho HTTP/3 sob netem. Os perfis degradados da submatriz (Seção 5.6) mostram IQR de p95 entre cerca de 11 ms e 19 ms. Isso exige cautela ao interpretar throughput e latência comparativos. O salto adicional e a configuração do proxy dialogam com estudos que integram HTTP/3 via proxy ou avaliam APIs REST em múltiplas versões do protocolo (MUÑOZ et al., 2025; TRIPATHI; MIRAZ; JOSHI, 2023)

Em relação ao desempenho em comparação ao HTTP/1.1 e HTTP/2 (QP3), com dez réplicas por célula, os dados coletados não permitem afirmar, de forma conclusiva, se HTTP/3 supera ou não HTTP/1.1 e HTTP/2 no arranjo experimental. Três fatores impedem essa conclusão: HTTP/1.1 foi medido sem passar pelo NGINX, enquanto HTTP/2 e HTTP/3 compartilham terminação TLS/QUIC na borda; o cliente HTTP/3 difere do k6; e o p95 de HTTP/3 partiu de centenas de milissegundos mesmo em clean, enquanto sua taxa de requisições não cresceu com a concorrência, o que sugere que o gerador de carga pesa mais que o efeito de rede simulado. Entre HTTP/2 e HTTP/3 no mesmo proxy, HTTP/2 apresentou menores p95 e maiores taxas de requisição nas condições testadas, mas essa diferença permanece sujeita às ameaças à validade (Seção 5.6). Os achados são compatíveis com a literatura que reporta ganhos dependentes de contexto, sem implicar vantagem universal (ZAHRA; ARDIANSYAH, 2026; PERNA et al., 2022; TRIPATHI; MIRAZ; JOSHI, 2023; ZHANG et al., 2025).

Já no que tange às estratégias de integração em ambientes consolidados (QP4), a estratégia adotada proxy reverso que termina HTTP/3 na borda e encaminha requisições à API Node.js - mostrou-se funcional para ambientes em que o \textit{runtime} permanece em HTTP/1.1 ou HTTP/2, sem exigir mudanças no código da aplicação. A configuração anuncia HTTP/3 via Alt-Svc e, em princípio, permite \textit{fallback} para versões anteriores, como discutido na fundamentação; os testes, porém, forçaram cada protocolo no cliente, de modo que a negociação e o \textit{fallback} não foram avaliados experimentalmente. Essa abordagem coincide com análises comparativas de HTTP/3 e HTTP/2 na presença de proxy, úteis para equipes que não alteram o núcleo da aplicação (MUÑOZ et al., 2025).

\section{Contribuições}

As três contribuições anunciadas na Introdução - revisão rápida da literatura, artefato reprodutível e evidências empíricas - foram entregues. Os resultados acrescentam a elas dois aprendizados metodológicos. O primeiro é a constatação de que, no arranjo testado, o gerador de carga QUIC limita a medida de HTTP/3 a ponto de mascarar o efeito do perfil de rede, o que torna a homogeneidade das ferramentas de carga um requisito metodológico, e não um detalhe. O segundo é a evidência de que a assimetria de topologia (HTTP/1.1 direto versus HTTP/2 e HTTP/3 via proxy) precisa ser controlada no desenho de qualquer comparação entre versões do HTTP em APIs Node.js. O código-fonte e os roteiros operacionais estão publicados em repositório aberto (\url{https://github.com/mendeslian/QUICBench-Node}), permitindo replicação e extensão das matrizes experimentais.

\section{Limitações}

A limitação principal é a restrição do desenho experimental: a matriz ideal testou variação de payload e concorrência apenas em rede clean, enquanto a submatriz degradada se ateve a GET /bench/ping com 50 usuários virtuais em quatro perfis de rede. As duas campanhas não se cruzam. Além disso, HTTP/1.1 foi medido com acesso direto à API, enquanto HTTP/2 e HTTP/3 passaram pelo NGINX; esse isolamento de topologia faz com que toda comparação entre HTTP/1.1 e os demais confunda efeito de protocolo com efeito de proxy. O cliente QUIC que mediu HTTP/3 difere do k6 em paralelismo efetivo, o que distorce a latência e a vazão relativos ao HTTP/2 no mesmo proxy. Por fim, os testes rodaram na mesma estação local em Docker, sem replicação em outro hardware. As respostas às questões de pesquisa, em especial QP3, refletem esse arranjo específico, não a viabilidade de HTTP/3 em produção.

A revisão rápida da literatura foi feita por um único aluno, sem segundo revisor na triagem, na leitura integral, na avaliação de qualidade e na extração. Cartaxo, Pinto e Soares (2020) admitem esse arranjo em revisão rápida, com risco de viés de seleção.

\section{Considerações finais}

Os resultados responderam às quatro questões de pesquisa em um arranjo específico de laboratório, não em ambiente de produção. O código-fonte aberto e os roteiros disponibilizados permitem que outros pesquisadores validem ou estendam o experimento, o que é o principal ganho desta rodada. O que fica pendente é medir HTTP/1.1 pelo mesmo NGINX dos demais protocolos, equiparar os geradores de carga e sair de localhost todos passos que o artefato já foi projetado para suportar.


